\documentclass[11pt]{article}
% Draft methods/results, styled for a Remote Sensing of Environment submission.
% Ready for reformatting into elsarticle.cls at submission time; body text and
% citation commands (natbib \citep/\citet) are format-agnostic.
\usepackage[margin=2.5cm]{geometry}
\usepackage{natbib}
\usepackage{booktabs}
\usepackage{amsmath}
\bibliographystyle{apalike}

\title{Predicting coverage-standardized plant diversity from raw Landsat phenology
curves across a unified Chilean vegetation plot network}
\author{}
\date{}

\begin{document}
\maketitle

\section{Methods}

\subsection{Study data and diversity facets}

We combined two vegetation plot networks spanning central and southern Chile.
Parcelas-CL \citep{cerdaparedes2026parcelascl} is a community-sourced database of 1082
plots between 30\textdegree{}S and 38\textdegree{}S. Living Trees is a forest inventory
of 2020 plots with a substantially wider geographic extent. We merged both networks into
a unified pool of 3102 plots, matching plot identifiers across networks by exact
coordinates.

We derived three diversity facets following the methodology of
\citet{perezgiraldo2026environmental}, applied to the subset of the unified pool with
genuine individual-level abundance counts (479 Parcelas-CL plots recorded under the
Abundance stratum, plus individual tree counts from Living Trees). We estimated the
local contribution to beta diversity (LCBD) using the quantitative Sørensen
decomposition of Podani and Schmera, implemented in the \emph{adespatial} package
(\texttt{beta.div.comp(coef="S", quant=TRUE)} followed by
\texttt{LCBD.comp(sqrt.D=TRUE)}; \citealp{legendre2013beta}). This facet
(\emph{lcbd\_count\_sorensen}) covers 2499 of the 3102 plots. We estimated
coverage-standardized phylogenetic and taxonomic diversity with iNEXT.3D
\citep{chao2021measuring}, following \citet{perezgiraldo2026environmental}:
\texttt{estimate3D(base="coverage", level=NULL, nboot=1)} for Hill orders $q=0,1,2$.
Phylogenetic diversity (\emph{pd\_inext\_q0/q1/q2}) additionally requires restricting
the species pool to a 610-tip phylogeny built with V.PhyloMaker2
\citep{jinqian2022vphylomaker2} and covers 888 plots; taxonomic diversity
(\emph{td\_inext\_q0/q1/q2}), without that restriction, covers 895. The
coverage-based estimator requires a minimum of five observed species per sampling unit,
which accounts for the lower coverage of both facets relative to
\emph{lcbd\_count\_sorensen}. We reindexed all three facets onto the full 3102-plot pool,
leaving missing values where the estimator did not meet its species floor, so that no
plot row is ever dropped.

\subsection{Remote sensing predictors and topography}

We extracted Landsat surface reflectance time series (NDVI, EVI, kNDVI, NBR and SAVI)
for the centre pixel of a 5$\times$5 Landsat window at each plot. We interpolated each
raw series onto a regular grid of 100 steps with a five-observation moving-average
smoothing window, without fitting any parametric composite curve; a preliminary
comparison against a composite phenological curve showed the raw interpolated curve
consistently outperforming the composite, so we used the raw curve throughout.

In addition to the phenological curve, we built a context block from
digital-elevation-model-derived topography (elevation, slope, northness, eastness,
heat-load index, topographic position index, ruggedness and curvature, each summarised
as the window mean and standard deviation) and plot area. We extracted Living Trees
topography independently and merged it with the Parcelas-CL topography through an exact
coordinate match, removing a previous median-imputation step that had affected all 2020
Living Trees plots.

\subsection{Cross-validation}

We evaluated every model with \texttt{kfold5\_block20\_unified}: five folds grouped by
20$\times$20\,km geographic blocks over a Chile-centred equal-area projection, following
the spatial-blocking principle for avoiding leakage from spatial autocorrelation
\citep{roberts2017crossvalidation}. This blocking scheme groups the 3102 plots into 610
distinct blocks.

\subsection{Predictive models}

We compared four model families under the same design (raw kNDVI curve, centre pixel,
topography+area context, \texttt{kfold5\_block20\_unified}):

\begin{itemize}
  \item Random Forest \citep{breiman2001random} on the concatenated curve and context.
  \item A 1-D convolutional network over the 100-step curve.
  \item A 2-D convolutional network (separable architecture, $\sim$15k parameters) over
    the curve reshaped as an image via a serpentine transform.
  \item The same 2-D network, initialised from a Masked Autoencoder
    \citep{he2022masked} pretrained on kNDVI curves from the non-unified pool.
\end{itemize}

We trained each model with three to five random seeds and report the pooled
out-of-sample $R^2$ across plots.

\section{Results}

\subsection{Vegetation index selection}

We swept all five vegetation indices under the same design (raw curve, centre pixel,
real topography, \texttt{kfold5\_block20\_unified}). We averaged $R^2$ only over the
three facets with usable signal ($R^2 > 0.4$: \emph{lcbd\_count\_sorensen},
\emph{pd\_inext\_q0}, \emph{td\_inext\_q0}); the remaining facets show no signal under
any index, and averaging them with the informative facets would have masked the only
signal present.

\begin{table}[h]
\centering
\begin{tabular}{lrrrr}
\toprule
Index & \emph{lcbd\_count\_sorensen} & \emph{pd\_inext\_q0} & \emph{td\_inext\_q0} & Mean ($R^2>0.4$) \\
\midrule
kNDVI & 0.434 & 0.608 & 0.772 & \textbf{0.6047} \\
EVI   & 0.426 & 0.612 & 0.775 & 0.6043 \\
NDVI  & 0.433 & 0.608 & 0.767 & 0.6027 \\
SAVI  & 0.424 & 0.615 & 0.763 & 0.6007 \\
NBR   & 0.435 & 0.604 & 0.725 & 0.5880 \\
\bottomrule
\end{tabular}
\caption{$R^2$ by vegetation index (raw curve, centre pixel, \texttt{kfold5\_block20\_unified}).}
\end{table}

kNDVI achieved the highest mean $R^2$, by a narrow margin over EVI (0.6047 versus
0.6043); NBR ranked lowest, driven by its weaker performance on \emph{td\_inext\_q0}. We
adopted kNDVI as the default index for all subsequent analyses.

\subsection{Effect of real Living Trees topography}

Replacing the median-imputed topography of Living Trees plots with measured topography
increased \emph{lcbd\_count\_sorensen} by 0.02--0.04 across all five indices, with no
notable effect on \emph{pd\_inext\_q0} or \emph{td\_inext\_q0}. We retained real
topography as the default configuration.

\subsection{Model comparison}

With kNDVI fixed, we compared the four model families under
\texttt{kfold5\_block20\_unified}:

\begin{table}[h]
\centering
\begin{tabular}{lrrrr}
\toprule
Facet & 2D-CNN & RF & 1D-CNN & 2D-CNN + MAE \\
\midrule
\emph{lcbd\_count\_sorensen} & 0.434 & \textbf{0.447} & 0.420 & 0.437 \\
\emph{pd\_inext\_q0}         & 0.608 & 0.582 & 0.612 & \textbf{0.613} \\
\emph{pd\_inext\_q1}         & 0.011 & 0.004 & \textbf{0.049} & -0.010 \\
\emph{pd\_inext\_q2}         & 0.065 & 0.054 & \textbf{0.091} & 0.040 \\
\emph{td\_inext\_q0}         & 0.772 & 0.666 & 0.780 & \textbf{0.786} \\
\emph{td\_inext\_q1}         & -0.178 & \textbf{0.042} & -0.144 & -0.181 \\
\emph{td\_inext\_q2}         & -0.079 & -0.073 & -0.066 & -0.090 \\
\bottomrule
\end{tabular}
\caption{Model family comparison ($R^2$, kNDVI, raw curve, centre pixel,
\texttt{kfold5\_block20\_unified}).}
\end{table}

The masked-autoencoder-initialised network achieved the highest absolute $R^2$ on
\emph{pd\_inext\_q0} (0.613) and \emph{td\_inext\_q0} (0.786), without transferring that
gain to \emph{lcbd\_count\_sorensen} or to $q1$/$q2$. The 1-D network matched or exceeded
the 2-D network on richness despite having fewer parameters (14{,}535 versus 15{,}175)
and was the only family with usable signal on \emph{pd\_inext\_q1/q2} (0.049 and 0.091).
Random Forest trained in 114\,s against 8--10\,min for each neural network and won on
\emph{lcbd\_count\_sorensen} (0.447), though it lost phylogenetic and taxonomic richness
to all three neural architectures.

\subsection{Synthesis}

Coverage-standardized richness (\emph{pd\_inext\_q0}, \emph{td\_inext\_q0}) is the only
facet family robustly predictable under this design, with $R^2$ between 0.58 and 0.79
across all four model families and all five indices tested. \emph{lcbd\_count\_sorensen}
offers moderate, stable signal (0.39--0.45). Abundance- and dominance-weighted facets
(\emph{pd\_inext\_q1/q2}, \emph{td\_inext\_q1/q2}) show no usable signal under any model,
index or architecture tested -- a limitation of the signal available in those specific
facets, not of model architecture.

\bibliography{refs}

\end{document}
